\section{The matroid \texorpdfstring{$M(\tilde G)$}{M(G~)}, deficiency, and \texorpdfstring{$k$}{k}-dof graphs}
\label{sec:molecular-deficiency}

This chapter is the \textbf{Phase~19} dep-graph: the combinatorial
substrate of the molecular-conjecture program ---
Katoh--Tanigawa~\cite[\S2.5, \S3]{katohTanigawa2011}, the third
stratum. Where \cref{sec:molecular-rigidity-matrix} built the
\emph{analytic} side (the rigidity matrix $R(G,p)$, its rank, and the
three rank lemmas), this chapter builds the \emph{matroidal} side: the
matroid $M(\tilde G)$, the $D$-deficiency $\operatorname{def}(\tilde
G)$, the $k$-degree-of-freedom hierarchy, rigid subgraphs, and the
bridge $\operatorname{def}(\tilde G) = \operatorname{corank} M(\tilde
G)$ that ties the two sides together. With that bridge in hand, the
reconciliation node deferred from Phase~18
(\cref{prop:rigidity-matrix-prop11}) closes. Throughout
$D = \binom{d+1}{2} = \delta$ is the body-bar dimension and $\tilde G$
abbreviates the multiplied graph $(D-1)\cdot G$ of
\cref{sec:body-hinge}, whose edge-fibers $\tilde e$ are the $D-1$
parallel copies of each edge $e \in E(G)$.

\medskip
\noindent\emph{Status (forward mode).} All nodes in this chapter are
red --- this dep-graph is the Phase~19 to-do list. The Lean lands in a
new file \texttt{Molecular/Deficiency.lean}; nodes flip to green as the
phase proceeds, leaf-first.

\subsection{The matroid \texorpdfstring{$M(\tilde G)$}{M(G~)} at the boundary regime}
\label{sec:molecular-deficiency-matroid}

\begin{definition}[The matroid $M(\tilde G)$]
  \label{def:matroid-MG}
  \lean{Graph.matroidMG, Graph.matroidMG_indep_iff}
  \leanok
  \uses{thm:tay-witness}
  Let $D = \binom{d+1}{2}$. The matroid $M(\tilde G)$ is the
  \emph{$(D,D)$-count matroid} of the multiplied graph
  $\tilde G = (D-1)\cdot G$ --- equivalently, the $D$-fold union of the
  cycle (graphic) matroid of $\tilde G$, restricted to $E(\tilde G)$.
  Its independent sets are the $(D,D)$-sparse edge subsets of
  $\tilde G$; by Tutte--Nash-Williams (\cref{thm:tay-witness} and the
  tree-packing chapter) these are exactly the edge subsets that
  decompose into $D$ edge-disjoint forests.
\end{definition}
\begin{proof}
  \leanok
  This is the \textbf{boundary regime} $\ell = 2k = D$ of the
  count-matroid family. It is \emph{not} covered by the matroidal-regime
  construction of \cref{sec:count-matroid} (built for $\ell < 2k$);
  $M(\tilde G)$ is instead the $D$-fold graphic-matroid union of
  \cref{sec:body-bar,sec:body-hinge}, whose independence
  characterization is Tutte--Nash-Williams tree-packing. The
  first formalization step of Phase~19 confirms this boundary regime is
  clean before the rest of the chapter relies on it.
\end{proof}

\begin{definition}[$D$-deficiency]
  \label{def:D-deficiency}
  \lean{Graph.partitionDef, Graph.deficiency}
  \leanok
  \uses{def:matroid-MG}
  For a partition $P$ of $V(G)$ into nonempty parts, write $d_G(P)$ for
  the number of edges of $G$ joining distinct parts of $P$. The
  \emph{$D$-deficiency} of $P$ is
  \[
    \operatorname{def}_{\tilde G}(P) \;=\; D\,(|P| - 1) \;-\;
      (D-1)\,d_G(P),
  \]
  and the \emph{deficiency} of $\tilde G$ is
  $\operatorname{def}(\tilde G) = \max_{P} \operatorname{def}_{\tilde
  G}(P)$, the maximum over all partitions. The deficiency measures how
  far $\tilde G$ is from being independently realizable as a rigid
  body-hinge framework. A partition is encoded by a labeling
  $f \colon V \to V$ whose fibers are the parts; $|P|$ is the number of
  distinct labels on $V(G)$ and $d_G(P)$ the number of edges with
  differently-labeled endpoints. The quantity
  $\operatorname{def}_{\tilde G}(P)$ is genuinely signed (a fine
  partition crossing many edges is negative), so it is $\mathbb Z$-valued;
  the maximum is $\ge 0$, witnessed by the trivial one-part partition.
\end{definition}

\begin{definition}[$k$-dof and minimal $k$-dof graphs]
  \label{def:k-dof}
  \lean{Graph.IsKDof, Graph.IsMinimalKDof, Graph.edgeFiber}
  \leanok
  \uses{def:D-deficiency, def:matroid-MG}
  A multigraph $G$ is a \emph{$k$-dof-graph} when
  $\operatorname{def}(\tilde G) = k$; the $0$-dof case is exactly the
  body-hinge-rigid case (\cref{thm:body-hinge-tay}). $G$ is a
  \emph{minimal $k$-dof-graph} when it is $k$-dof and every base $B$ of
  $M(\tilde G)$ meets every edge-fiber $\tilde e$ (the $D-1$ parallel
  copies of $e \in E(G)$) --- i.e.\ no edge of $G$ can be deleted
  without changing the deficiency. Minimal $k$-dof-graphs are the
  objects the combinatorial induction (\cref{sec:molecular-deficiency},
  Phase~20's Theorem~4.9) reduces to the two-vertex double edge.
\end{definition}

\subsection{Rigid subgraphs and structural lemmas}
\label{sec:molecular-deficiency-rigid}

\begin{definition}[Rigid and proper rigid subgraphs]
  \label{def:rigid-subgraph}
  \uses{def:k-dof}
  A subgraph $H \subseteq G$ is \emph{rigid} when it is $0$-dof
  ($\operatorname{def}(\tilde H) = 0$), i.e.\ $\tilde H$ packs $D$
  edge-disjoint spanning trees; it is a \emph{proper} rigid subgraph
  when $\emptyset \ne V(H) \subsetneq V(G)$. A \emph{circuit} of
  $M(\tilde G)$ is a minimal dependent edge set. Two-edge-connectivity
  of $G$ is the structural hypothesis Case~III of the algebraic
  induction works under.
\end{definition}

\begin{lemma}[Two-edge-connectivity; KT Lemma 3.1]
  \label{lem:two-edge-conn}
  \uses{def:k-dof}
  A minimal $k$-dof-graph $G$ with $|V| \ge 2$ and $k \ge 0$ is
  $2$-edge-connected.
\end{lemma}
\begin{proof}
  A bridge would split off a part contributing to the deficiency that
  could be realized independently, contradicting minimality; this is
  Katoh--Tanigawa~\cite[Lemma 3.1]{katohTanigawa2011}.
\end{proof}

\begin{lemma}[Subgraph minimality via restriction; KT Lemma 3.3]
  \label{lem:subgraph-minimality}
  \uses{def:k-dof, def:matroid-MG}
  A rigid subgraph $H$ of a minimal $k$-dof-graph $G$ is itself a
  minimal $0$-dof-graph, via the restriction $M(\tilde G)|E(\tilde H) =
  M(\tilde H)$.
\end{lemma}
\begin{proof}
  Matroid restriction commutes with the union construction, so the base
  / fiber-meeting property descends to $H$; Katoh--Tanigawa~\cite[Lemma
  3.3]{katohTanigawa2011}.
\end{proof}

\begin{lemma}[A circuit yields a rigid subgraph; KT Lemma 3.4]
  \label{lem:circuit-rigid}
  \uses{def:rigid-subgraph, def:matroid-MG}
  The edge set of a circuit of $M(\tilde G)$ spans a rigid subgraph of
  $G$.
\end{lemma}
\begin{proof}
  A circuit is one element short of $D$ edge-disjoint spanning trees on
  its vertex span, hence is $0$-dof there; Katoh--Tanigawa~\cite[Lemma
  3.4]{katohTanigawa2011}.
\end{proof}

\subsection{The deficiency \texorpdfstring{$=$}{=} corank bridge}
\label{sec:molecular-deficiency-bridge}

\begin{theorem}[$\operatorname{def}(\tilde G) = \operatorname{corank} M(\tilde G)$; Jackson--Jord\'an]
  \label{thm:def-eq-corank}
  \uses{def:D-deficiency, def:matroid-MG}
  For any base $B$ of $M(\tilde G)$,
  \[
    |B| + \operatorname{def}(\tilde G) \;=\; D\,(|V| - 1),
  \]
  i.e.\ $\operatorname{def}(\tilde G) = \operatorname{corank} M(\tilde
  G)$ inside the rank-$D(|V|-1)$ ambient. This is the project framing of
  the generic-rank identity of Jackson--Jord\'an~\cite[Theorem
  6.1]{jacksonJordan2009}, whose Corollary~6.2 reads off rigidity
  ($\operatorname{def} = 0$) as $\tilde G$ packing $D$ edge-disjoint
  spanning trees.
\end{theorem}
\begin{proof}
  The rank of the $D$-fold graphic-matroid union is the partition
  formula of Edmonds (\cref{sec:matroid-union}); the corank against the
  ambient $D(|V|-1)$ is exactly the deficiency maximised over
  partitions. The conjecture (Theorem~5.6) needs only the upper-bound
  half of the generic-rank identity, which the body-hinge reduction of
  \cref{thm:body-hinge-tay} already supplies; the prove-vs-hypothesize
  boundary for the full Jackson--Jord\'an~\cite{jacksonJordan2009}
  identity is decided when this node lands.
\end{proof}

This bridge closes the reconciliation node deferred from Phase~18:
\cref{prop:rigidity-matrix-prop11} in
\cref{sec:molecular-rigidity-matrix}, which equates the honest rank
form $\operatorname{rank} R(G,p) = D(|V|-1)$ with the standard-basis
existence form (\cref{thm:body-hinge-tay}) of Phase~16, routes through
$\operatorname{def}(\tilde G) = \operatorname{corank} M(\tilde G)$ ---
once \cref{thm:def-eq-corank} is green, the body-hinge analytic and
combinatorial sides agree.
